%----------------------------------------------------------------
%% Template para preparação de trabalho de conclusão de 
%% curso   (TCC)   do   Departamento   de   Economia   da 
%% Universidade      Estadual     de      Londrina       (UEL)
%
%% Autores: 
%% Prof Ph.D. Joanna G. Alexopoulos
%% Prof. Dr. Marcelo S. Bego 
%
%% Data da última atualização: 22/02/2020
%----------------------------------------------------------------
%%
%% Customizações do abnTeX2 (http://abnTeX2.googlecode.com) 
%%
%% Este trabalho pode ser distribuído e/ou modificado sob as
%% condições da Licença Pública do LaTeX Project, versão 1.3.
%% A versão desta licença pode ser encontrada em:
%%   http://www.latex-project.org/lppl.txt
%% e faz parte da versão LaTeX 01/12/2005.
%%
%% Este trabalho tem como status de manutenção LPPL como "maintained".
%% 
%% Os atuais responsáveis por esse trabalho são: 
%% Prof Ph.D. Joanna G. Alexopoulos
%% Prof. Dr. Marcelo S. Bego 
%%
%% Mais informações sobre abnTeX2 estão disponíveis em: 
%%  https://github.com/abntex/abntex2
%%
%----------------------------------------------------------------
%% Sobre a classe abntex2.cls:
%% abntex2.cls, v-1.9.5 laurocesar
%% Copyright 2012-2015 by abnTeX2 group at https://www.abntex.net.br/ 
%%
%----------------------------------------------------------------

\documentclass[12pt,oneside,a4paper,chapter=TITLE,english,brazil,sumario=abnt-6027-2012]{abntex2}

% Dark mode
\usepackage{xcolor}
\pagecolor[rgb]{0.1,0.1,0.15} %black
\color[rgb]{0.5,0.5,0.5} %grey

% %%%%%%%%%%%%%%%%
% PACOTES
% %%%%%%%%%%%%%%%%

\usepackage{lmodern}			
\usepackage[T1]{fontenc}		
\usepackage[utf8]{inputenc}		
\usepackage{indentfirst}		
\usepackage[dvipsnames]{xcolor}				
\usepackage{graphicx}			
\usepackage{microtype} 
\usepackage{multicol}
\usepackage{multirow}
\usepackage{float}	
\usepackage{lipsum}				
\usepackage{tikz}
\usepackage{ragged2e} 
\usepackage[brazilian,hyperpageref]{backref}	
\usepackage[alf]{abntex2cite}
\usepackage{blindtext} 
\usepackage{mathptmx}
\usepackage[labelfont=bf]{caption}
\captionsetup{labelfont=bf}
\usepackage[singlelinecheck=false]{caption}
\usepackage{amsmath}

% %%%%%%%%%%%%%%%%
%MARGENS
%As margens são configuradas conforme a NBR 14724:2011
%As margens devem ser: para o anverso, esquerda e superior de 3 cm e direita e inferior de 2 cm; para o verso, direita e superior de 3 cm e esquerda e inferior de 2 cm. 

% %%%%%%%%%%%%%%%%
%FONTES 
% %%%%%%%%%%%%%%%%

% %%%%%%%%%%%%%%%%
%TIPO DE FONTE
\renewcommand{\sfdefault}{\rmdefault} % Fonte: Times New Roman

% %%%%%%%%%%%%%%%%
% TAMANHO E ESTILO DAS FONTES DOS TÍTULOS E SUBSEÇÕES 
\renewcommand{\ABNTEXchapterfontsize}{\large \bfseries}%14pt

\renewcommand{\ABNTEXsectionfont}{\bfseries}%Subseções em negrito.

% %%%%%%%%%%%%%%%%
% ESPAÇAMENTOS
% %%%%%%%%%%%%%%%%

% TAMANHO DO PARÁGRAFO 
\setlength{\parindent}{3cm} %Primeira linha

%ESPAÇO ENTRE LINHAS
%O padrão é definido como \OnehalfSpacing (1.5)


% %%%%%%%%%%%%%%%%
%INFORMAÇÕES DA CAPA, FOLHA DE ROSTO E FOLHA DE APROVAÇÃO 
% %%%%%%%%%%%%%%%%

% %%%%%%%%%%%%%%%%
%NOME DO TRABALHO 
\titulo{\bfseries ANÁLISE DA DINÂMICA INFLACIONÁRIA SOB O REGIME DE METAS DE INFLAÇÃO NO BRASIL ENTRE 19XX e 20XX}


% %%%%%%%%%%%%%%%%
%NOME DO AUTOR 
\autor{LUCAS BONI DOS ANJOS AMARAL ALVARENGA}


% %%%%%%%%%%%%%%%%
%NOME DO ORIENTADOR
\orientador{CARLOS EDUARDO CALDARELLI}


% %%%%%%%%%%%%%%%%
\local{LONDRINA, PARANÁ}
\data{\the\year}
\tipotrabalho{monografia}
\instituicao{Universidade Estadual de Londrina (UEL)}
\preambulo{ Trabalho de Conclusão de Curso apresentado ao Departamento de Economia da Universidade  Estadual de Londrina.}
% %%%%%%%%%%%%%%%%


% %%%%%%%%%%%%%%%%
% CAPA UEL
\renewcommand{\imprimircapa}{
	\begin{capa}
		\center
		\begin{figure}
		\includegraphics[height=4.5cm,width=14cm]{Logo_UEL}
		\end{figure}

		\begin{tikzpicture}
		\fill[ForestGreen!90] (17.1,0.6) rectangle (1,1);
		\end{tikzpicture}
		\large 
		
		{\bfseries CENTRO DE ESTUDOS SOCIAIS APLICADOS\\
		DEPARTAMENTO DE ECONOMIA\\
		CURSO DE CIÊNCIAS ECONÔMICAS\\}
		
		
		\vspace{5cm}
		
		{\ABNTEXchapterfont\textsc{\Large\imprimirtitulo}}
		
		\vfill
		
		\begin{flushright}
			{\ABNTEXchapterfont\textsc{\Large \imprimirautor}}
			
		\end{flushright}
		
		\begin{tikzpicture}
		\fill[ForestGreen!90] (17.1,0.6) rectangle (1,1);
		\end{tikzpicture}
		
		{\imprimirlocal}
		
		{\imprimirdata}
		
	\end{capa}
}

% %%%%%%%%%%%%%%%%
% FOLHA DE ROSTO
\makeatletter

\renewcommand{\folhaderostocontent}{
	\begin{center}
      \large 	
	
	  {\ABNTEXchapterfont\textsc{\Large \imprimirautor}}
		
		\vspace{6cm}
		
		{\ABNTEXchapterfont\textsc{\Large \imprimirtitulo}}
		
		\vspace{2cm}
		
		\begin{flushright}
			\begin{minipage}{8cm}
				\SingleSpacing
				Trabalho de Conclusão de Curso apresentado ao Departamento de Economia da Universidade Estadual de Londrina.
				
				Orientador: Prof(a). \imprimirorientador
			\end{minipage}%
		\end{flushright}
		
		\vfill
		
		
		{\large\imprimirlocal}
		
		{\large\imprimirdata}
	\end{center}
}
\makeatother


% %%%%%%%%%%%%%%%%
% FOLHA DE APROVAÇÃO 
\makeatother

\newcommand{\folhaDeaprovacao}{
	\begin{center}
		
		\begin{folhadeaprovacao} 
			\begin{center} 
				
					
				
			{\ABNTEXchapterfont\textsc{\large \imprimirautor}}
				
				\vspace{2cm}
				
			{\ABNTEXchapterfont\textsc{\large \imprimirtitulo}}
				
				\vspace{1cm}		
				\begin{flushright}
					\begin{minipage}{10cm}
						\SingleSpacing
						Trabalho de Conclusão de Curso apresentado ao Departamento de Economia da Universidade Estadual de Londrina.
					\end{minipage}%
				\end{flushright}	
				
				\vspace{2cm}
				
				\begin{flushright}
					\begin{minipage}{10cm}		
						\centering
					{\bfseries COMISSÃO EXAMINADORA}
					\end{minipage}
				\end{flushright}	
				
				\vspace{1cm}
				
				\begin{flushright}
					\begin{minipage}{10cm}	
						\centering
						\hrule \hspace{0.2cm}
						
						Orientador(a): Prof(a). \imprimirorientador  
						
						\imprimirinstituicao
					\end{minipage}%
				\end{flushright}
				
				\vspace{1cm}
				
				\begin{flushright}
					\begin{minipage}{10cm}	
						\centering
						\hrule \hspace{0.2cm}
						
						Prof(a). NOME BANCA 1
						
						\imprimirinstituicao
					\end{minipage}%
				\end{flushright}		
				
				
				\vspace{1cm}
				
				\begin{flushright}
					\begin{minipage}{10cm}	
						\centering
						\hrule \hspace{0.2cm}
						
						Prof(a). NOME BANCA 2 
						
						\imprimirinstituicao
					\end{minipage}%
				\end{flushright}
				
				\vfill 
				
				\begin{flushright}
					\begin{minipage}{6cm}	
				\centering
				Londrina, XX de XX de \imprimirdata.
					\end{minipage}
			\end{flushright}
		
			\end{center} 
		\end{folhadeaprovacao}
	\end{center}
}
\makeatother


% %%%%%%%%%%%%%%%%
%REGULAMENTO DO TRABALHO DE CONCLUSÃO DO CURSO DE ECONOMIA (DISCIPLINAS MONOGRAFIA I E II – 6TCC404 E 6TCC405)
% %%%%%%%%%%%%%%%%

%ART.16 A estrutura da monografia compõe-se de:

%I. Capa;
%II. Folha de rosto;
%III. Elementos pré-textuais;
%IV. Resumo e “Abstract” (ou versão em Espanhol ou ainda em Francês):
%V. Sumário;
%VI. Introdução;
%VII. Desenvolvimento, contendo Referencial Teórico ou Econométrico ou Revisão de Literatura ou Revisão Bibliográfica ou Modelo Teórico; e nos casos de estudos quantitativos Metodologia e Fonte e tratamento de dados;
%VIII. Resultados e discussão;
%IX. Considerações finais ou Conclusão;
%X. Referências;
%XI. Anexos e Apêndices, quando for o caso
%%%%%%%%%%%%%%%%%%%%%%%%%%%%



% INÍCIO DO DOCUMENTO
\begin{document}

% Retira espaço extra obsoleto entre as frases.
\frenchspacing 


% %%%%%%%%%%%%%%%%
% CAPA
\imprimircapa

% %%%%%%%%%%%%%%%%
%FOLHA DE ROSTO
\folhaderostocontent


% %%%%%%%%%%%%%%%%
%FIXA CATALOGRÁFICA 
%Incluir fica catalográfica
%\pagebreak


% %%%%%%%%%%%%%%%%
%FOLHA DE APROVAÇÃO 
\folhaDeaprovacao


% %%%%%%%%%%%%%%%%
%DEDICATÓRIA
\begin{dedicatoria}
	\vspace*{\fill}
	\noindent
	Dedico esse trabalho aos meus pais, Bruno e Silvia, por todo amor, apoio e incentivo, que foram fundamentais para a realização de mais uma conquista. Sua dedicação e sacrifícios me proporcionaram a oportunidade de seguir em frente e realizar meus objetivos. 
	\vspace*{\fill}
\end{dedicatoria}


% %%%%%%%%%%%%%%%%
%AGRADECIMENTOS
\begin{agradecimentos}
	\noindent
	Agradeço aos professores e colegas de curso, que contribuíram com conhecimentos, experiências e incentivos, tornando a jornada acadêmica mais rica e gratificante. Aos amigos, que estiveram presentes nos momentos de lazer e de estudo, proporcionando equilíbrio e motivação durante o percurso.	À Universidade Estadual de Londrina por oferecer um ambiente de aprendizado e crescimento. A todos aqueles que, direta ou indiretamente, colaboraram para a conclusão deste trabalho.
\end{agradecimentos}


% %%%%%%%%%%%%%%%%
%RESUMO
\noindent
ALVARENGA, Lucas. {\bfseries Análise da Dinâmica Inflacionária Sob o Regime de Metas de Inflação no Brasil Entre 19XX e 20XX}, \imprimirdata. <FOLHAS> f. Monografia (Curso de Ciências Econômicas). Centro de Estudos Sociais Aplicados, Universidade Estadual de Londrina, Londrina, 2024.

\setlength{\absparsep}{18pt} % ajusta o espaçamento dos parágrafos do resumo
\begin{resumo}

	Este trabalho tem por objetivo verificar o impacto da inflação de oferta e da inflação de demanda sobre a economia brasileira e avaliar os efeitos exercidos por estas sobre a eficácia do regime de metas de inflação (RMI) no Brasil. Será examinado se as respostas do Banco Central são eficientes e sob quais óticas a política monetária deve ser definida a fim de atender os objetivos propostos pelo RMI. Por meio deste trabalho, busca-se atingir um melhor entendimento dos mecanismos que regem a dinâmica inflacionária no Brasil e como é moldada a resposta à inflação pela autoridade monetária brasileira.
	
	\textbf{Palavras-chave}: inflação; oferta; demanda
\end{resumo}


% %%%%%%%%%%%%%%%%
%ABSTRACT
\noindent
ALVARENGA, Lucas. {\bfseries Análise da Dinâmica Inflacionária Sob o Regime de Metas de Inflação no Brasil Entre 19XX e 20XX}, \imprimirdata. <FOLHAS> f. Monografia (Curso de Ciências Econômicas). Centro de Estudos Sociais Aplicados, Universidade Estadual de Londrina, Londrina, 2024.

\begin{resumo}[Abstract]
	\begin{otherlanguage*}{english}
		
		Write your abstract here... 
		
		\vspace{\onelineskip}
		
		\noindent 
		\textbf{Keywords}: 
	\end{otherlanguage*}
\end{resumo}
\pagebreak


% %%%%%%%%%%%%%%%%
%LISTA DE FIGURAS
\pdfbookmark[0]{\listfigurename}{lof}
\listoffigures*
\cleardoublepage


% %%%%%%%%%%%%%%%%
%LISTA DE TABELAS
\pdfbookmark[0]{\listtablename}{lot}
\listoftables*
\cleardoublepage


% %%%%%%%%%%%%%%%%
%LISTA DE ABREVIATURAS
\begin{siglas}
	\item[UEL] Universidade Estadual de Londrina. 
	\item[ABNT] Associação Brasileira de Normas Técnicas.
	\item[RMI] Regime de metas de inflação
	\item[IPCA] Índice Nacional de Preços ao Consumidor Amplo
\end{siglas}
\pagebreak


% %%%%%%%%%%%%%%%%
%LISTA DE SíMBOLOS
\begin{simbolos}
	\item[Depeco] Departamento de Economia
\end{simbolos}
\pagebreak


% %%%%%%%%%%%%%%%%
% SUMÁRIO
\pdfbookmark[0]{\contentsname}{toc}
\tableofcontents*
\cleardoublepage


% %%%%%%%%%%%%%%%%
%INÍCIO DO TEXTO
\textual % indica o início do texto para a numeração das páginas
\pagestyle{simple}
\aliaspagestyle{chapter}{simple}

\chapter{Introdução}

Diversos países adotam, atualmente, regimes de metas de inflação. O modelo surgiu em 1990 na Nova Zelândia e chegou ao Brasil em 1999, quando foi proposta uma série de diretrizes que têm por objetivo atingir uma meta preestabelecida de crescimento de um índice de inflação, nomeadamente o Índice Nacional de Preços ao Consumidor Amplo (IPCA) \cite{sicsu2002teoriarmi}. Não há consenso econômico entre economistas acerca da teoria econômica que é mais adequada para sustentar o RMI. Tem-se discordância, por exemplo, se de fato existe uma taxa natural de desemprego, abaixo de cuja desencadear-se-ia um processo inflacionário.

\chapter{Tipos de Inflação}

\section{Inflação de Oferta}

\subsection{Subseção}

\subsubsection{Subsubseção}

\section{Inflação de Demanda}

\section{Inflação Estrutural}

\section{Inflação Inercial}

\chapter{Regime de Metas de Inflação}

\section{RMI No Brasil}

\section{Política Monetária}

\chapter{Metodologia}

\section{Modelo}

\section{Dados}

\chapter{Considerações Finais}



%%%%%%%%%%%%%
%BIBLIOGRAFIA
\bibliographystyle{abntex2-alf}
\bibliography{ref}



%%%%%%%%%%%%%%%%%
%APÊNCICES
\apendices

\chapter{Anexos e Apêndices, quando for o caso.}

\chapter{Cronograma}

\begin{table}[H]
	\centering
	\caption{Cronograma}
	\begin{tabular}{|c|c|c|c|c|c|c|c|c|c|c|c|}
		\toprule
		Atividades\textbackslash{}Data & Jan & Fev & Mar & Abr & Mai & Jun & Jul & Ago & Set & Out & Nov \\
		\midrule
		Atividade 1                    &     &     &     &     &     &     &     &     &     &     &     \\
		\midrule
		Atividade 2                    &     &     &     &     &     &     &     &     &     &     &     \\
		\midrule
		                               &     &     &     &     &     &     &     &     &     &     &     \\
		\midrule
		                               &     &     &     &     &     &     &     &     &     &     &     \\
		\midrule
		                               &     &     &     &     &     &     &     &     &     &     &     \\
		\midrule
		                               &     &     &     &     &     &     &     &     &     &     &     \\
		\midrule
		                               &     &     &     &     &     &     &     &     &     &     &     \\
		\midrule
		                               &     &     &     &     &     &     &     &     &     &     &     \\
		\midrule
		                               &     &     &     &     &     &     &     &     &     &     &     \\
		\midrule
		                               &     &     &     &     &     &     &     &     &     &     &     \\
		\midrule
		                               &     &     &     &     &     &     &     &     &     &     &     \\
		\midrule
		                               &     &     &     &     &     &     &     &     &     &     &     \\
		\bottomrule
	\end{tabular}%
	\label{tab:addlabel}%
\end{table}%


\chapter{Dicas}

\section {Figuras e Tabelas}

\subsection{Figuras}

Para adicionar uma figura, crie o ambiente de figura e utilize \verb|\includegraphics| para anexar a figura no texto. Com esse comando você pode definir a largura e altura da figura no texto. 

Exemplo: 

\vspace{0.2cm}

\verb|\begin{figure}[H]|

\verb|\caption{Logo da UEL}|

\verb|\includegraphics[width=7cm,height=2cm]{Logo_Uel}\\|

\verb|{\footnotesize Fonte: ... }|

\verb|\end{figure}|

\vspace{0.5cm}

Compilando esses comandos chegamos ao seguinte resultado: 


\begin{figure}[H]
	\caption{Logo da UEL}
	\includegraphics[width=7cm,height=2cm]{Logo_Uel}\\
	{\footnotesize Fonte: ... }
\end{figure}

{\bfseries Importante:} As figuras devem estar na mesma pasta que o TCC está salvo.	

\subsection{Tabelas}

Escrever tabelas no {\LaTeX} sempre foi um ato de coragem dado o tempo necessário para completar a tarefa. 

Para escrever uma tabela você precisa primeiro criar o ambiente tabular, definir quantas colunas indicando seu posicionamento (\verb|l = esquerda; c=centro; r=direita|) e adicionar os dados separados por \verb|&| e no final de cada linha indicar que ela terminou escrevendo \verb|\\| (quebra de linha). Linhas verticais podem ser geradas adicionando o caracter | e linhas horizontais o comando \verb|\hline|.

\vspace{0.2cm}

Exemplo:

\verb|\begin{table}[H]|

\verb|\caption{Nossa Tabela}|

\verb|\begin{tabular}{ c c }|

\verb|\hline|

\verb|1 & 2 \\|

\verb|4 & 5 \\|

\verb|\hline|

\verb|\end{tabular}\\|

\verb|{\footnotesize Nota: ... }|

\verb|\end{table}|

\vspace{0.2cm}

Compilando esses comando chegamos ao seguinte resultado: 

\vspace{0.2cm}

\begin{table}[H]
	\caption{Nossa Tabela}
	\begin{tabular}{ c  c }
		\hline
		1 & 2 \\
		4 & 5 \\
		7 & 8 \\
		\hline
	\end{tabular}\\
	{\footnotesize Nota: ... }
\end{table}
\vspace{0.5cm}

Entretanto uma maneira mais fácil é utilizar um conversor.  No nosso caso vamos utilizar um conversor de EXCEL par {\LaTeX}. Baixe o conversor Excel2LaTeX no seguinte link: {$\href{https://ctan.org/tex-archive/support/excel2latex?lang=en}{https://ctan.org/tex-archive/support/excel2latex?lang=en}$} e adicione o conversor aos \textit{Add-ins} do Excel. 


\section{Ambiente Matemático}

Quando estamos no meio do texto podemos criar um ambiente matemático utilizando cifrões. Por exemplo, \verb|$Y_t=K_t^\alpha L_t^{1-\alpha}$| sairá no texto como $Y_t=K_t^\alpha L_t^{1-\alpha}$. Observe que os caracteres \verb|_| gera subscrito e \verb|^| gera o sobrescrito. Entretanto, para escrevermos equações que mereçam destaque, devemos criar o ambiente equation, o qual numerará automaticamente as equações.  

\vspace{0.2cm}

Exemplo:

\vspace{0.2cm}

\verb|\begin{equation}|

\verb|Y_t=K_t^\alpha L_t^{1-\alpha}|

\verb|\end{equation}|

\vspace{0.2cm}

Compilando esses comando chegamos a seguinte equação

\vspace{0.2cm}

\begin{equation}\label{equ1}
	Y_t=K_t^\alpha L_t^{1-\alpha}
\end{equation}

Algumas operações que normalmente usamos são frações, somatórias, limites, derivas e integrais:

\begin{itemize}
	\item Fração = \verb|\frac{numredaor}{denominador}| = $\frac{Numerador}{Denominador}$
	      
	\item Somatória = \verb|\sum_{mínimo}^{máximo} | = $\sum_{min}^{max}$
	      
	\item Derivada = \verb|\partial| = $\partial$
	      
	\item Integral, \verb|\int_{mínimo}^{máximo}| = $\int_{min}^{max}$
	      
\end{itemize}

No caso das matrizes precisamos criar o ambiente de matrizes, bmatrix e escreve-la de forma parecida como escrevemos uma tabela.

\vspace{0.2cm}
Exemplo: 
\vspace{0.2cm}

\verb|\[|

M=

\verb|\begin{bmatrix}|

\verb|1 & 2  \\|

\verb|3 & 4 |

\verb|\end{bmatrix}|

\verb|\]|

\vspace{0.4cm}

Nossa Matriz:

\vspace{0.2cm}
\[
	M=
	\begin{bmatrix}
		1 & 2   \\
		3 & 4 
	\end{bmatrix}
\]




\section{Citações e Referências}\label{citaebib}

{\tiny }

Para fazer citações e a seção de referências é necessário utilizar um pacote de apoio as várias opções. No nosso caso estamos utilizando o pacote \verb|\usepackage[alf]{abntex2cite}|, o qual traz as normas da ABNT.

\subsection{Arquivo .bib}

Para gerar sua bibliografia é necessário criar um aquivo que contêm a sintaxe padrão da extensão .bib das referências que serão utilizadas. Cada referência deve apresentar uma lista de informações que serão utilizadas na geração das referências. 

O exemplo abaixo mostra a entrada no aquivo .bib de duas teses de doutorado. Cada entrada é definida como uma tese, @phdthesis, que será chamada pelo {\LaTeX} como alexopoulos2012three e a outra bego2017three. Nas referências serão utilizadas as informações do título, nome do autor, ano e instituição.  

\vspace{1cm}
{
\noindent
@phdthesis{alexopoulos2012three,\\
title={Three essays on inequality},\\
author={Alexopoulos, Joanna},\\
year={2012},\\
school={University of Illinois at Urbana-Champaign}\\
}
}

\vspace{1cm}
{
\noindent
@phdthesis{bego2017three,\\
title={Three essays on agricultural markets},\\
author={Bego, Marcelo da Silva},\\
year={2017}\\
}
}


Você não precisa gerar a entrada de cada artigo, pois o Google Scholar oferece esse formato. Em cite no Google Scholar procure por BibTeX. 

Nomeando e salvando o aquivo que contêm as referências de ref.bib (ou qualquer nome que você deseja) na mesma pasta que o TCC está salvo é possível adicionar a citação no texto e o \LaTeX \hspace{0.1cm} gerará a referência automaticamente. 

	{{\bfseries Importante:} Para compilar o aquivo ref.bib use F8}. Assim use F8 para compilar a entrada de novas referências e F5  para compilar o texto com a nova entrada de referência e atualizar o pdf.  

\subsection{Citações}

Para fazer a citação com \verb|\usepackage[alf]{abntex2cite}| temos as seguintes opções: \verb|\citet|,  \verb|\citeonline|.

Exemplos: 

\verb|\cite{alexopoulos2012three}| = (ALEXOPOULOS, 2012)

\verb|\citeonline{alexopoulos2012three}| = Alexopoulos (2012)

\verb|\cite{bego2017three}| = (BEGO, 2017)

\verb|\citeonline{bego2017three}| = Bego (2017)

\subsection{Referências}

A bibliografia é gerada e atualizada automaticamente cada vez que adicionamos ou retiramos uma citação no corpo do texto. Para gerar a seção referências utiliza-se o seguintes comandos:  

\vspace{1cm}
{
	\noindent
	\verb|\bibliographystyle{abntex2-alf}%Definição do padrão ANBT| \\
	\verb|\bibliography{ref}%chamando o arquivo .bib com as referências | \\
}

Como resultado as referências bibliográficas do texto serão apresentadas da seguinte forma: 

\vspace{1cm}
\begin{center}
	\large \bfseries REFERÊNCIAS
\end{center}

{
\noindent
ALEXOPOULOS, J. Three essays on inequality. Tese (Doutorado)—University of Illinois at
Urbana-Champaign, 2012.

\noindent
BEGO, M. d. S. Three essays on agricultural markets. Tese (Doutorado), 2017.
}


\end{document}
